\section{Test Case 24: long run to a stationary bead count}

Test Case 24 is a long production run whose purpose is to reach a
statistically stationary active-bead count.  Once the jet has bridged the
16 cm nozzle-to-collector gap with a developed bending instability, insertion
at the nozzle and removal at the collector are expected to balance, so that
the bead count oscillates around a mean instead of drifting.  Unlike Test
Cases 21--23 it is not designed to finish within a smoke-test timeframe: it
integrates \texttt{final time 0.5d0} seconds with \texttt{timestep 5.d-9}
seconds, about 100 million steps.

The case is assembled from three sources.  Its electrostatic parameters are
JETSPIN's canonical set, \texttt{density charge 44000},
\texttt{collector distance 16}, and \texttt{external potential 30.02076857},
shared with Examples 3, 4 and 5 and with Test Cases 9--20.  Rheology,
air drag, insertion and removal are those of Test Case 23: Maxwell rheology,
stochastic Platen air drag, nozzle insertion, and collector removal, with
its anchor spacing, \texttt{dynamic refinement anchor 0.10} cm.
Initialization and print directives are those of Example 3, so the jet grows
from a single nozzle bead (\texttt{points 1}) rather than from the
pre-extended 400-element mesh used to validate the refinement algorithm.  This
last point is what distinguishes the case: every other anchored
stochastic-Platen refinement test (Test Cases 21--23) starts from an already
extended jet, so the young-jet growth regime of this configuration had not
been exercised.  Example 5 also grows from one bead with refinement, but
with RK4, multiple-step Coulomb sums, and no anchors.

Yarin evaporation, which Test Case 23 uses, is not enabled here, and this is a
definition rather than a temporary state; the evaporation directives of the
input are read but unused.  The evaporative configuration is Test Case 25
(Sec.~\ref{sec:Test-Case-25}), whose input differs from this one by two
lines, \texttt{evaporation yes} and \texttt{evaporation polymer frac 0.50d0}.
With the 6\% polymer fraction of Ref.~\cite{yarin2001bending} the evaporative
configuration stopped on a numerical instability near step 1.07 million in
runs of a build of 29 September 2026, because the dried jet keeps its charge
while losing fifteen times its mass and cross-section; Test Case 25 carries
that analysis.

The dynamic-refinement cadence is taken from Example 5,
\texttt{dynamic refinement threshold 0.4} cm and
\texttt{dynamic refinement every 1.d-3} s, rather than Test Case 23's tighter
\texttt{0.10} cm and \texttt{1.d-5} s.  This choice is load-bearing rather
than cosmetic.  With the tighter cadence the jet accumulates dozens of
accepted refinement events within fewer than 200,000 steps, each one
re-fitting the cross-section radius from the previous event's own output.  The
result is a nonphysical thinning compounding far beyond the
order-of-magnitude reduction that electrospinning genuinely produces,
collapsing bead mass toward zero until the resulting acceleration overflows
the stress equation.  The advisory \texttt{warning(108)} is raised whenever a
refinement threshold is set below twenty times the base discretization
resolution.  It is informational and not enforced: Example 5 and Test Cases
24 and 25 sit exactly at that ratio (0.4 cm over a resolution of 0.02 cm)
and raise no warning, while Test Cases 21--23, at five times the resolution,
are below it and remain valid short-window references.  Refinement also
starts late: \texttt{dynamic refinement start 2.5d-4}, a directive that
neither Example 5 nor Test Case 23 sets, allows no refinement during the
first $2.5\times10^{-4}$ s (50,000 steps).

Since 5 October 2026 the case follows Test Case 25: CPU and GPU builds read
the same pre-generated Gaussian pool (Sec.~\ref{sec:Noise-pool}), and the GPU
runs the device-resident path once the jet holds 100 beads.  The reference
run uses the NVFORTRAN 24.3 CPU build and is truncated at 5,000,000 steps, 5\%
of the target duration, by \texttt{final time 2.5d-2} and
\texttt{print time 1.d-4} in a copy of the input; it closes correctly.  The
first bead reaches the
collector at step 2,541,113.  Between 4 and 5 million steps the collector
velocity is 1963 cm/s, the off-axis distance of the leading bead 3.8 cm, the
cone angle 26.6 degrees, the active-bead count between 471 and 553 around a
mean of 511, and the path length 211 cm, so the balanced insertion and removal
regime is reached; four more CPU seeds give mean active counts between 511.3
and 512.0.  One NVIDIA
A30 runs the same 5 million steps in 596 s, with the process bound to the
NUMA node of its GPU: it reproduces the CPU run
exactly until step 1,379,105, where it engages the device-resident path,
follows the CPU trajectory afterwards (with seed 317 the active-bead count
first differs at step 4.28 million; until 7 October 2026 the bead inserted at
the array reallocation of step 1,963,389 came one step late on the GPU, and
the counts separated at 2.66 million steps), and five seeds give between 511
and 512 active beads and a path length of 211 cm over the same window.  The
A30 step then costs between about 105 and 145~$\mu$s from 100 to 600 beads,
against 4 to 5 ms for the CPU build with about 500 beads; the standard GPU
build of the morning of 5 October 2026, which kept this case on the host
with only the Coulomb sums on the GPU, needed 4546 s.  With NVHPC 25.5 the
run takes 630 s, against 597 s with NVHPC 24.3, the reference compiler, in
the same session, and gives the same insertion and removal totals and 511.2
active beads between 4 and 5 million steps, as NVHPC 24.3 does.  On one A30
the full 100 million steps took 14,906 s, about four hours, with a
device-resident build of 5 October 2026, whose output equals the current build's up to the
reallocation at step 1,963,389, where that build inserted one step late:
from 10 to 100 million steps the jet keeps 517 active beads on average
(474 to 560), a path length of 214 cm, an off-axis distance of 3.8 cm and a
cone angle of 26.9 degrees, about 1\% above the values between 4 and 5
million steps, where the jet is still settling.  These long-run statistics
rest on periodically repeating noise: with the default pool of $10^{8}$
values (Sec.~\ref{sec:Noise-pool}) the noise repeats every
$10^{8}/(6\,n_{\mathrm{active}})$ steps, about 32,000 steps at 517 beads, so
the 100 million steps go through the pool about 3,100 times.

Track the printed \texttt{n} column over time to observe the stationary
oscillation.  The input is stored in
\texttt{examples/input-24/input.dat}.  The investigation into refinement
robustness under repeated remeshing, which produced the cadence used here, is
recorded in \texttt{docs/refinement-robustness-investigation.md}.
